%======================================================================
%  Hardware platform: SD650 V3 and Intel Xeon CPU Max
%======================================================================

\section{Xeon Max and the SD650 V3}
\label{sec:xeon-max-platform}

\subsection{The unit of analysis}

The Lenovo ThinkSystem SD650~V3 is a dense, direct-water-cooled server
tray rather than a single compute node.  One tray contains two
independent dual-socket nodes, and six trays fit in a 6U DW612S
enclosure~\cite{lp1603}.  This distinction matters when reporting
quantities: the CRESCO8 allocation is a \emph{node}; memory locality is
defined first at the \emph{socket}; and the scheduler may allocate
individual cores within the node.  \tabref{tab:hardware-hierarchy} sets
out the four levels and what each one is good for.

\begin{table}[htbp]
\centering
\small
\tabsetup
\begin{tabularx}{\textwidth}{P{3cm}YY}
\headrow
\headcell{Level} & \headcell{Composition} & \headcell{Relevance} \\[3pt]
\tabheadrule
\tabstripeRow
DW612S enclosure
  & Six SD650~V3 trays (12 independent nodes)
  & Shared power and cooling infrastructure; not an allocation unit. \\[3pt]

\tabrowsep
\tabplainRow
SD650~V3 tray
  & Two independent dual-socket nodes
  & High-density packaging; the two nodes do not form one shared-memory
    system. \\[3pt]

\tabrowsep
\tabstripeRow
Compute node
  & Two Xeon processors connected by Intel UPI
  & Scheduler allocation and the largest cache-coherent NUMA system used
    by the single-node experiments. \\[3pt]

\tabrowsep
\tabplainRow
Xeon Max socket
  & CPU cores, 64\,GB on-package HBM2e and eight DDR5 channels
  & Primary locality domain and the cleanest unit for comparing HBM with
    DDR5. \\[3pt]
\tabbottomrule
\end{tabularx}
\caption{Hardware hierarchy and the properties relevant to this study.
Product capability is separated from the installed CRESCO8
configuration described in \secref{sec:cresco8}
~\cite{lp1603,cresco8doc}.}
\label{tab:hardware-hierarchy}
\end{table}

The product guide covers several processor, memory, storage and network
options.  Most of those options are irrelevant to the thesis.  The
important configuration is the Xeon CPU Max 9480 installed on CRESCO8:
56 physical cores per socket, 64\,GB of HBM2e, eight DDR5-4800 channels,
and four UPI links~\cite{intel9480,lp1603}.  A two-socket node therefore
provides 112 physical cores and 128\,GB of nominal HBM capacity.  The
node also contains DDR5, but its installed capacity and channel
population are properties of CRESCO8, not of the product family.

\begin{figure}[htbp]
\centering
\begingroup
\definecolor{socketink}{HTML}{263445}
\definecolor{socketblue}{HTML}{365D85}
\colorlet{socketaccent}{sapred}
\definecolor{socketborder}{HTML}{CBD5DF}
\resizebox{\textwidth}{!}{%
\begin{tikzpicture}[
  font=\small\sffamily, text=socketink,
  card/.style={line width=0.65pt, rounded corners=4pt, align=center,
               inner xsep=7pt, inner ysep=8pt},
  socket/.style={card, draw=socketborder, fill=socketink!3,
                 minimum width=4.6cm, minimum height=4.0cm},
  hbm/.style={card, draw=socketaccent, fill=socketaccent!9,
              minimum width=4.6cm, minimum height=1.1cm},
  ddr/.style={card, draw=socketblue, fill=socketblue!9,
              minimum width=4.6cm, minimum height=1.1cm},
  core/.style={draw=socketblue!60, fill=socketblue!14,
               line width=0.35pt, rounded corners=0.7pt,
               minimum width=0.24cm, minimum height=0.24cm,
               inner sep=0pt, outer sep=0pt},
  note/.style={font=\footnotesize\sffamily, text=socketink!80,
               align=center},
  link/.style={{Stealth[length=5pt,width=4pt]}-{Stealth[length=5pt,width=4pt]},
               line width=0.8pt, draw=socketink!75}
]
  \node[socket] (s0) at (-3.8,0) {};
  \node[socket] (s1) at (3.8,0) {};
  \foreach \i/\x in {0/-3.8,1/3.8} {
    \node[font=\small\sffamily\bfseries] at (\x,1.52) {Socket \i};
    \node[note] at (\x,1.12) {56 physical cores};
    % Eight columns and seven rows: one tile for each physical core.
    % The grid conveys multiplicity, not the processor's physical floorplan.
    \foreach \row in {0,...,6} {
      \foreach \col in {0,...,7} {
        \node[core] at ({\x-1.05+0.30*\col},{0.65-0.30*\row}) {};
      }
    }
    \node[note] at (\x,-1.62) {AMX instructions\quad |\quad LLC};
  }
  \node[hbm, below=0.5cm of s0] (h0)
    {\textbf{64\,GB HBM2e}\\{\footnotesize high bandwidth}};
  \node[hbm, below=0.5cm of s1] (h1)
    {\textbf{64\,GB HBM2e}\\{\footnotesize high bandwidth}};
  \node[ddr, above=0.5cm of s0] (d0)
    {\textbf{8 DDR5 channels}\\{\footnotesize high capacity}};
  \node[ddr, above=0.5cm of s1] (d1)
    {\textbf{8 DDR5 channels}\\{\footnotesize high capacity}};

  \draw[link] (s0.east) -- node[above, note]{UPI} (s1.west);
  \draw[link, draw=socketaccent] (s0.south) -- (h0.north);
  \draw[link, draw=socketaccent] (s1.south) -- (h1.north);
  \draw[link, draw=socketblue] (s0.north) -- (d0.south);
  \draw[link, draw=socketblue] (s1.north) -- (d1.south);

  \node[draw=socketborder, densely dashed, line width=0.5pt,
        rounded corners=6pt, inner sep=9pt,
        fit=(d0)(d1)(h0)(h1)] (compute) {};
  \node[note, below=0.18cm of compute]
    {One cache-coherent compute node};
\end{tikzpicture}%
}
\endgroup
\caption{Conceptual locality map of a CRESCO8 Xeon Max node.  Each
socket has local HBM and DDR5; access to the other socket traverses UPI.
Each small tile represents one physical core.  The diagram shows locality
and core count, not physical core placement or package routing.
Author's synthesis based on the platform and cluster documentation
~\cite{lp1603,cresco8doc}.}
\label{fig:node-locality}
\end{figure}


\subsection{What the hardware looks like}
\label{sec:hardware-photographs}

The internal layout in \figref{fig:sd650-internal} makes the locality
model concrete without turning the chapter into a product tour.  It is
the one product photograph retained in the main text because the
placement of processors, DIMMs and cold plates directly supports the
experimental argument.

\begin{figure}[htbp]
\centering
% Crop the border and drop shadow embedded in the source image.
\includegraphics[width=0.94\textwidth,trim=6bp 27bp 27bp 6bp,clip]{sd650v3-internal-view.png}
\caption{Inside the tray.  Each of the two servers carries two
processors, and each processor is flanked by eight memory DIMMs.  The
copper cold plates and the loop that joins them replace the fans
entirely.  The high-bandwidth memory that this thesis is about is
nowhere visible: it sits inside the processor package, beneath the cold
plate, which is why its capacity is fixed and why it cannot be
configured, replaced or upgraded.  Reproduced from the Lenovo
ThinkSystem SD650~V3 product guide~\cite{lp1603}.}
\label{fig:sd650-internal}
\end{figure}

The layout of one server, two processors with their own DIMMs on either side, is
the physical form of the locality argument made throughout this section:
a thread running on the left-hand processor reaches the DIMMs beside it
quickly and the DIMMs on the right slowly, and no amount of software can
change the physical distance.  Direct water cooling supports sustained
operation, but it does not guarantee an invariant clock; frequency,
temperature and throttling still have to be observed during measurement.

\subsection{Compute capability, and why it is not the binding constraint}

Xeon Max is a Sapphire Rapids processor with conventional x86 cores,
AVX-512, and Intel Advanced Matrix Extensions (AMX).  AMX adds per-core
tile registers and matrix-multiply instructions for BF16 and INT8
operands~\cite{intelmax,na2024llmcpu}, and on well-shaped
matrix--matrix products it is a large win, which is why prompt
processing benefits from it and single-token generation largely does
not.

The detail that matters for a quantized workload is what AMX does not
provide.  There is no native INT4 dot-product instruction, so a 4-bit
GGUF model cannot be fed to the matrix engine in the format it is stored
in.  \figref{fig:amx-path} shows the consequence: every weight must
pass through an unpack and dequantize step before it reaches a tile, and
that step sits precisely between the memory tier and the matrix unit.
The arrangement has an awkward property for this thesis.  Lower-bit
formats reduce the bytes moved, which is what makes them attractive on a
bandwidth-limited machine, but they add work on the path from memory to
the matrix engine, and they push the operand format further from
anything the engine can consume directly.  Whether a given model
benefits from AMX at all therefore depends on the runtime, the storage
format, the packing scheme and the matrix dimensions, none of which are
properties of the processor.

\begin{figure}[htbp]
\centering
\resizebox{\textwidth}{!}{%
\begin{tikzpicture}[
  font=\small,
  mem/.style={draw=hbmcol!85!black, fill=hbmcol!14, rounded corners=3pt,
              minimum width=2.3cm, minimum height=1.5cm, align=center},
  stage/.style={draw=black!65, fill=neutralcol, rounded corners=3pt,
                minimum width=2.7cm, minimum height=1.5cm, align=center},
  eng/.style={draw=black!70, fill=white, rounded corners=3pt,
              minimum width=2.4cm, minimum height=1.5cm, align=center},
  flow/.style={-{Stealth}, thick, black!70},
  edgelbl/.style={font=\scriptsize, align=center, text=black!70}
]
  \node[mem]   (m) at (0,0)    {HBM2e\\\emph{or}\\DDR5};
  \node[stage] (u) at (4.3,0)  {unpack \&\\dequantize\\{\footnotesize AVX-512}};
  \node[stage] (t) at (8.6,0)  {tile registers\\{\footnotesize BF16 / INT8}};
  \node[eng]   (x) at (12.1,0) {AMX\\\textsc{tmul}};

  \draw[flow] (m) -- node[above=1pt, edgelbl] {4-bit}
                    node[below=1pt, edgelbl] {packed} (u);
  \draw[flow] (u) -- node[above=1pt, edgelbl] {widened} (t);
  \draw[flow] (t) -- (x);

  \draw[dashed, thick, black!45, -{Stealth}]
    (m.north) .. controls +(2.4,1.6) and +(-2.8,1.6) .. (x.north);
  \node[edgelbl, text=black!55, fill=white, inner sep=2pt] at (6.05,1.62)
    {no native INT4 path};
\end{tikzpicture}%
}
\caption{Path taken by a quantized weight from a memory tier to the
matrix engine on Xeon Max.  The unpack stage is unavoidable for 4-bit
formats and lies between the two things this thesis compares.  At batch
one the right-hand operand is a vector, so the tile is largely empty and
the engine runs far below its rated throughput; the step is paced
instead by the arrival of weights from the left.  Author's synthesis
from the Xeon Max product documentation and the AMX description in Na et
al.~\cite{intelmax,na2024llmcpu}.}
\label{fig:amx-path}
\end{figure}

\textbf{Peak instruction throughput is not an application result.}  Small
matrices, batch-one matrix--vector operations, packing overhead and
memory stalls can all leave the matrix units idle, and the advertised
maximum turbo frequency is a single-core limit rather than the frequency
of a socket running sustained AMX kernels.  Frequency, package power and
throttling state have to be observed during the experiments rather than
read off a specification table.


\subsection{The memory hierarchy}
\label{sec:memory-hierarchy}

HBM and DDR5 trade different resources.  HBM provides much greater
aggregate bandwidth and 64\,GB of nominal capacity per socket.  DDR5
provides substantially more capacity and can have lower unloaded access
latency.  HBM is therefore not simply ``faster memory''.  It is most valuable
under two conditions that have to hold \emph{together}: the workload
issues enough concurrent requests to approach its bandwidth, \emph{and}
the bandwidth-critical working set is placed locally.  Either one alone
is not enough.

Published STREAM results illustrate both the advantage and the need for
local measurement.  Na et al.\ report 588\,GB/s for HBM and
233.8\,GB/s for DDR5 on a Xeon Max 9468 socket
~\cite{na2024llmcpu}.  Ibeid et al.\ report 680--840\,GB/s for HBM and
235--245\,GB/s for DDR5 on a related Xeon Max 9470C
~\cite{ibeid2025aurora}.  These sustained values are 2.5--3.4 times
apart, but the HBM results themselves differ by more than 15\%.  SKU,
clustering mode, thread placement and benchmark configuration are
plausible causes.  None of these numbers should be used as the bandwidth
of CRESCO8 before the Max 9480 nodes are measured.

Three paths must be distinguished in every result:

\begin{itemize}
  \item \textbf{local HBM or local DDR5}, reached from cores on the same
        socket;
  \item \textbf{remote memory}, reached through UPI from the other
        socket; and
  \item \textbf{mixed placement}, in which pages, threads or model
        objects span several locality domains.
\end{itemize}

\figref{fig:node-locality} draws these three paths.  An apparently
node-wide 128\,GB HBM allocation is consequently two 64\,GB local
tiers.  A workload uses the aggregate capacity and
bandwidth only if data and computation are balanced across sockets.

\begin{table}[htbp]
\centering
\small
\tabsetup
\begin{tabularx}{\textwidth}{P{2.35cm}P{3.05cm}P{3.15cm}Y}
\headrow
\headcell{Level} & \headcell{Capacity} & \headcell{Placement decided by} &
\headcell{Why it matters here} \\[3pt]
\tabheadrule
\tabstripeRow
L1 data & 48\,KB per core & Hardware & Too small to hold any part of a
  weight tensor. \\[3pt]

\tabrowsep
\tabplainRow
L2 & 2\,MB per core & Hardware & Holds tiles and activations; invisible
  to a placement policy. \\[3pt]

\tabrowsep
\tabstripeRow
Last-level cache & 112.5\,MB, shared & Hardware replacement policy &
  The tier Zhang et al.\ target by rebuilding the
  runtime~\cite{zhang2026cacheresident}; not addressable directly. \\[3pt]

\tabrowsep
\tabplainRow
HBM2e & 64\,GB, four stacks & \textbf{Software}, in Flat mode &
  The bandwidth tier, and the object of this study. \\[3pt]

\tabrowsep
\tabstripeRow
DDR5 & Eight channels; installed capacity is a CRESCO8 property &
  \textbf{Software}, in Flat mode & The capacity tier and the control
  condition. \\[3pt]

\tabrowsep
\tabplainRow
Remote memory & The other socket's tiers & \textbf{Software}, through
  thread and page binding & Reached over UPI; a placement failure rather
  than a design choice. \\[3pt]
\tabbottomrule
\end{tabularx}
\caption{Storage levels visible to one Max 9480 socket, ordered by the
question that matters here: who decides what occupies them.  Only the
last three rows are under software control, and only in Flat mode.
Cache capacities follow the SKU specification and should be confirmed
with \texttt{lscpu} on the allocated node~\cite{intel9480,lp1603}.}
\label{tab:storage-levels}
\end{table}

\tabref{tab:storage-levels} also locates this thesis against the
closest competing work.
Zhang et al.\ obtain large gains by keeping weights resident in the
last-level cache~\cite{zhang2026cacheresident}, one row above the tier
studied here, but that row is not addressable, so their result required
rebuilding the runtime around it.  The HBM row can be
addressed directly, which is what makes the comparison in this thesis
possible without touching the runtime at all.

\subsection{Bandwidth is only reachable with enough concurrency}
\label{sec:littles-law}

A bandwidth number describes what the memory system can deliver when it
is kept busy, and nothing about whether a given program manages to keep
it busy.  The relationship is Little's law: the number of bytes a
program must have in flight to sustain a bandwidth \(B\) against a
latency \(L\) is

\begin{equation}
\label{eq:little}
  C = B \times L .
\end{equation}
\eqcaption{Little's law applied to a memory system.  It converts a bandwidth figure into a requirement on the program: this many bytes must be outstanding at every instant, or the bandwidth is unreachable no matter which tier holds the data.}

Substituting the published Xeon Max figures into \eqnref{eq:little}
makes the difference between the two tiers concrete.  For HBM at 840\,GB/s and 134.9\,ns, roughly
113\,KB must be outstanding at all times, which is about 1{,}770 cache
lines, or close to 32 lines per core across 56 cores.  For DDR5 at 245\,GB/s
and 112.9\,ns the requirement is roughly 28\,KB, about 430 lines, or
under 8 per core~\cite{ibeid2025aurora}.  The faster tier does not merely
offer \emph{more} bandwidth; it demands roughly \textbf{four times the
memory-level parallelism} before that bandwidth becomes reachable at
all.

That demand is not obviously satisfiable.  Thirty-two concurrent line
fills per core is the same order of magnitude as the outstanding-miss
resources a single core has available, so approaching HBM peak requires
almost every one of them to stay occupied, helped by hardware
prefetchers that work well on strided streams and poorly on dependent
ones.  The exact resource counts for this microarchitecture must be
confirmed rather than assumed, but the shape of the constraint does not
depend on the precise number.

This is the mechanism behind the most important negative result in the
tiered-memory literature.  Peng et al.\ found that HPC codes with
regular access patterns gained up to 3$\times$ from MCDRAM while codes
with irregular, latency-bound patterns could run \emph{slower} in the
high-bandwidth tier, and that adding hardware threads, which is to say adding concurrency,
recovered much of the loss~\cite{peng2017hybridmemory}.  Low-batch autoregressive decode has
exactly the awkward shape: the weight stream is regular and prefetchable,
but the dependent chain through each layer limits how far ahead the core
can run, and the attention path over a growing KV cache is less regular
still.

Two consequences follow for the experiments.  Thread count is not a
nuisance parameter to be fixed at ``all cores'' but a first-class
factor, because it directly sets available concurrency; and a
disappointing HBM result must be diagnosed by measuring achieved
bandwidth rather than reported as an absence of effect.  A run that
reaches 250\,GB/s on HBM has not established that HBM bandwidth was the
limiting resource.  Its traffic rate remains within the measured DDR5
envelope, so any performance difference cannot be attributed to
exploiting HBM's additional bandwidth alone.

\subsection{HBM memory modes}
\label{sec:hbm-modes}

The four 16\,GB HBM2e stacks are physically present in every Xeon Max
package, but software does not necessarily see them as 64\,GB of main
memory.  The firmware's \emph{memory hierarchy} setting determines the
role of HBM: it can be the only main memory, a separately addressable
tier alongside DDR5, or an invisible cache in front of DDR5.  These are
the three modes in \figref{fig:hbm-modes}.  The choice is a platform
configuration, normally fixed for the lifetime of a boot; it is not a
policy that an individual batch job can switch
~\cite{lp1738,intelmaxtuning}.

\begin{figure}[htbp]
\centering
\begingroup
\definecolor{modeink}{HTML}{263445}
\definecolor{modeblue}{HTML}{365D85}
\colorlet{modeaccent}{sapred}
\definecolor{modeborder}{HTML}{CBD5DF}
\resizebox{\textwidth}{!}{%
\begin{tikzpicture}[
  font=\small\sffamily, text=modeink,
  card/.style={line width=0.65pt, rounded corners=4pt, align=center,
               inner xsep=5pt, inner ysep=7pt},
  panel/.style={card, draw=modeborder, fill=white,
                minimum width=4.35cm, minimum height=5.15cm},
  cpu/.style={card, draw=modeborder, fill=modeink!3,
              minimum width=3.7cm, minimum height=0.85cm},
  hbm/.style={card, draw=modeaccent, fill=modeaccent!9,
              text width=3.3cm, minimum width=3.7cm,
              minimum height=1.2cm},
  ddr/.style={hbm, draw=modeblue, fill=modeblue!9},
  absent/.style={card, draw=modeborder, densely dashed,
                 text=modeink!55, fill=modeink!1,
                 text width=3.3cm, minimum width=3.7cm,
                 minimum height=1.2cm},
  summary/.style={card, draw=modeborder, fill=modeink!3,
                  font=\footnotesize\sffamily, text width=3.95cm,
                  minimum width=4.35cm, minimum height=1.65cm},
  link/.style={{Stealth[length=5pt,width=4pt]}-{Stealth[length=5pt,width=4pt]},
               line width=0.8pt, draw=modeink!75},
  note/.style={font=\footnotesize\sffamily, text=modeink!80, align=center}
]
  \foreach \x in {0,4.85,9.7} {
    \node[panel] at (\x,0.6) {};
  }
  \node[font=\small\sffamily\bfseries] at (0,3.95) {HBM-only};
  \node[font=\small\sffamily\bfseries] at (4.85,3.95) {Flat (1LM)};
  \node[font=\small\sffamily\bfseries] at (9.7,3.95) {Cache (2LM)};
  \node[note] at (0,3.5) {HBM is the only main memory};
  \node[note] at (4.85,3.5) {Two separate main-memory tiers};
  \node[note] at (9.7,3.5) {HBM is a cache for DDR5};

  \node[cpu] (c0) at (0,2.5) {\textbf{CPU cores}};
  \node[cpu] (c1) at (4.85,2.5) {\textbf{CPU cores}};
  \node[cpu] (c2) at (9.7,2.5) {\textbf{CPU cores}};

  % HBM-only: one connected, allocatable tier.  DDR5 is not used.
  \node[hbm] (h0) at (0,0.7)
    {\textbf{HBM}\\{\footnotesize 64\,GB main memory}};
  \node[absent] at (0,-1.0)
    {\textbf{DDR5}\\{\footnotesize not used / not visible}};
  \draw[link, draw=modeaccent] (c0.south) -- (h0.north);

  % Flat: independent, parallel paths to two allocatable NUMA regions.
  % Side-by-side placement avoids suggesting that DDR5 is behind HBM.
  \node[hbm, text width=1.55cm, minimum width=1.85cm] (h1) at (3.8,0.7)
    {\textbf{HBM}\\{\footnotesize 64\,GB}};
  \node[ddr, text width=1.55cm, minimum width=1.85cm] (d1) at (5.9,0.7)
    {\textbf{DDR5}\\{\footnotesize RAM}};
  \draw[link, draw=modeaccent] ($(c1.south)+(-1.05,0)$) -- (h1.north);
  \draw[link, draw=modeblue] ($(c1.south)+(1.05,0)$) -- (d1.north);
  \node[note, text width=3.7cm] at (4.85,-1.0)
    {Separate NUMA regions.\\Software chooses the tier\\for each allocation.};

  % Cache: HBM is a transparent memory-side cache, not a RAM allocation target.
  \node[hbm, densely dashed] (h2) at (9.7,0.7)
    {\textbf{HBM cache}\\{\footnotesize 64\,GB; not allocatable}};
  \node[ddr] (d2) at (9.7,-1.0)
    {\textbf{DDR5}\\{\footnotesize main memory}};
  \draw[link, draw=modeaccent] (c2.south) -- (h2.north)
    node[midway, right, note, xshift=3pt] {lookup};
  \draw[link, draw=modeblue] (h2.south) -- (d2.north)
    node[midway, right, note, xshift=3pt] {miss / fill};

  % Identical summary rows make visibility, capacity and control comparable.
  \node[summary] at (0,-3.0)
    {\textbf{OS sees:} HBM only\\
     \textbf{RAM:} 64\,GB\\
     \textbf{Placement:} HBM only};
  \node[summary] at (4.85,-3.0)
    {\textbf{OS sees:} HBM + DDR5\\
     \textbf{RAM:} 64\,GB + DDR5\\
     \textbf{Placement:} software};
  \node[summary] at (9.7,-3.0)
    {\textbf{OS sees:} DDR5 only\\
     \textbf{RAM:} DDR5 capacity\\
     \textbf{Placement:} hardware};
\end{tikzpicture}%
}
\endgroup
\caption{The three Xeon Max memory modes, shown for one socket.  Solid memory
cards denote allocatable RAM; the dashed HBM card in Cache mode is a
transparent, hardware-managed, direct-mapped cache.  Cache misses access DDR5,
and HBM cache capacity is not added to RAM.  Capacities are nominal; the paths
are conceptual and omit the CPU's internal caches.
Author's synthesis from the Xeon Max product
and tuning guides~\cite{lp1603,lp1738,intelmaxtuning}.}
\label{fig:hbm-modes}
\end{figure}

\begin{description}
  \item[HBM-only mode.] Firmware exposes HBM as the system's only usable
  main-memory tier; DDR5 is absent from the operating system's address
  space.  Lenovo documents this topology for a configuration without
  installed DDR5, while a site may implement an equivalent experimental
  state by administratively disabling populated DDR5.  The exact
  mechanism must therefore be recorded on the machine.  Applications
  require no HBM-specific placement mechanism, but the operating system,
  services and all jobs share the same 64\,GB per-socket capacity.  The
  usable application footprint is smaller than the nominal 64\,GB
  (128\,GB in a dual-socket node).  This is the simplest route to HBM
  bandwidth when the complete workload fits, but it is not the documented
  production CRESCO8 configuration~\cite{lp1738}.

  \item[Flat mode (1LM).] HBM and DDR5 are exposed as separate NUMA
  address spaces.  The capacity of both tiers is available, and software
  can put a bandwidth-critical object in HBM while leaving colder data in
  DDR5.  That control is also an obligation: the application or launch
  policy must bind threads and pages with \texttt{mbind}, libnuma or
  \texttt{numactl}.  First touch, page migration and file-backed mappings
  can still change actual residency, so requested placement must be
  verified rather than inferred from the command line.

  \item[Cache mode (2LM).] Software sees the DDR5 address space; the
  on-package HBM is an invisible, hardware-managed, direct-mapped
  memory-side cache.  Existing applications can therefore benefit
  without an HBM-aware allocator and can use the full DDR5 capacity, but
  software cannot reserve HBM for model weights or exclude cold objects.
  Conflict misses, warm state and a working set near the cache capacity
  can make performance workload-dependent
  ~\cite{lp1738,ibeid2025aurora}.
\end{description}

Flat mode requires two configuration stages on Linux.  Firmware must
first select \emph{Flat} under \emph{Memory Hierarchy}.  After boot, each
HBM DAX device must be changed from device-DAX to
\texttt{system-ram} with \texttt{daxctl}; on Lenovo's dual-socket example
this means two devices in Quadrant mode and eight in SNC4.  The operation
must be repeated after every boot and therefore belongs in the node's
administrator-managed startup configuration, not in an ordinary user
job.  Only after that step do the HBM regions appear as allocatable
memory-only NUMA nodes~\cite{lp1738}.

The right choice follows the working set and the software's control over
placement.  HBM-only removes tier selection when the complete workload
comfortably fits in HBM.  Flat mode supports larger workloads and gives
an HBM-aware program deterministic control over the hot set.  Cache mode
supports the same large capacity with a transparent policy when such
control is unavailable, at the cost of making HBM occupancy and
replacement implicit.  For this thesis Flat mode provides the strongest
experimental design because local HBM and local DDR5 can be compared on
the same socket without changing the processor or executable.  A
Flat-versus-Cache comparison answers a different question: explicit
placement versus a hardware cache policy.

\subsection{Clustering and NUMA}
\label{sec:hbm-clustering}

Memory mode is only the first configuration axis.  An independent
firmware setting controls how finely each socket is partitioned.  With
SNC disabled, which Lenovo calls \emph{Quadrant} or \emph{Quad} mode,
the socket is one compute locality domain even though the package still
contains four memory-controller quadrants.  SNC4 partitions the socket
into four sub-NUMA clusters, each with a subset of cores, last-level
cache and DDR channels.  When HBM is visible as main memory, each of the
four HBM stacks also becomes a 16\,GB locality domain; in Flat mode these
are additional memory-only nodes, whereas in Cache mode they remain
invisible.  SNC4 can increase local bandwidth and reduce contention for
a NUMA-aware program, but it also reduces the capacity of one local HBM
allocation from 64\,GB to 16\,GB and creates more ways to access memory
remotely~\cite{lp1738,intelmaxtuning}.

Crossing three memory modes with two clustering modes gives the six
configurations in \tabref{tab:hbm-six-configurations}.  The NUMA counts
are for a complete dual-socket node, because that is the unit allocated
on CRESCO8.  They describe the topology after any Flat-mode DAX devices
have been converted to system RAM.  The two axes must not be conflated:
Quadrant and SNC4 are not alternatives to HBM-only, Flat or Cache.
Rather, every memory mode can be paired with either clustering mode, so
an HBM-only experiment must still report whether it used HBM-only+Quad
or HBM-only+SNC4~\cite{lp1738}.

\begin{table}[htbp]
\centering
\small
\tabsetup
\begin{tabularx}{\textwidth}{P{1.8cm}P{1.55cm}P{4.35cm}Y}
\headrow
\headcell{Memory mode} & \headcell{Cluster} &
\headcell{Visible topology (two sockets)} &
\headcell{Operational consequence} \\[3pt]
\tabheadrule
\tabstripeRow
HBM-only & Quad
  & 2 HBM nodes: one 64\,GB node per socket
  & Bind work per socket; DDR5 is absent and all software consumes HBM. \\[3pt]

\tabrowsep
\tabplainRow
HBM-only & SNC4
  & 8 HBM nodes: four 16\,GB nodes per socket
  & Bind each worker and its allocation to the same quadrant. \\[3pt]

\tabrowsep
\tabstripeRow
Flat & Quad
  & 4 nodes: two DDR5 nodes with CPUs and two 64\,GB HBM memory-only nodes
  & Select the tier explicitly and pair each HBM node with cores on its
    local socket. \\[3pt]

\tabrowsep
\tabplainRow
Flat & SNC4
  & 16 nodes: eight DDR5 nodes with CPUs and eight 16\,GB HBM memory-only
    nodes
  & Select the tier and the matching quadrant; a misplaced page can be
    local to the socket but remote to the worker. \\[3pt]

\tabrowsep
\tabstripeRow
Cache & Quad
  & 2 visible DDR5 nodes, one per socket; HBM is invisible
  & Bind by socket, but HBM contents remain hardware-managed. \\[3pt]

\tabrowsep
\tabplainRow
Cache & SNC4
  & 8 visible DDR5 nodes, four per socket; HBM is invisible
  & Bind by quadrant; software still cannot choose what the HBM caches. \\[3pt]
\tabbottomrule
\end{tabularx}
\caption{The six Xeon Max HBM configurations and their OS-visible NUMA
topology on a dual-socket server.  ``Memory-only'' means that a NUMA node
contains capacity but no CPUs.  Node identifiers are assigned by the
running system and must be discovered rather than copied from the vendor
example.  Synthesised from Lenovo's mode and topology diagrams
~\cite{lp1738}.}
\label{tab:hbm-six-configurations}
\end{table}
\FloatBarrier

The two Flat-mode topologies are shown directly in Lenovo's hardware
diagrams in \figref{fig:lenovo-flat-topologies}.  In Quad mode, each
socket has one CPU-bearing DDR5 node and one 64\,GB HBM memory-only node.
In SNC4, each of those is split four ways.  The HBM node numbers shown in
the drawings are not architectural constants: on the illustrated boot,
DDR5 occupies nodes 0--7 and HBM nodes 8--15, but firmware, kernel and
hot-plug order can change numbering.  The stable relationship to recover
is which CPU-bearing node is closest to each memory-only HBM node.

Flat mode is particularly unforgiving because converting a DAX device
to system RAM creates the HBM node but does not assign it CPUs.  A
default allocation made by a thread on a CPU-bearing node can therefore
land in DDR5; selecting an HBM node without selecting nearby cores can
instead create remote traffic.  A valid experiment must inspect
\texttt{numactl --hardware}, derive the CPU--HBM distance pairing, bind
both CPU and memory, and then verify residency and traffic counters.

This distinction explains why a mode cannot be labelled universally
``best''.  Ibeid et al.\ obtain the highest STREAM bandwidth with SNC4
and explicit locality, whereas Na et al.\ observe worse inference
performance and more remote traffic under SNC4
~\cite{ibeid2025aurora,na2024llmcpu}.  The result depends on the
software's placement policy, not only on the processor setting.

\begin{figure}[htbp]
\centering
% Crop the outer border and drop shadow, preserving the NUMA boundaries.
\includegraphics[width=0.97\textwidth,trim=4.8bp 21.6bp 21.6bp 4.8bp,clip]{lenovo-lp1738-flat-quad.png}

\vspace{6pt}
\includegraphics[width=0.90\textwidth,trim=4.8bp 21.6bp 21.6bp 4.8bp,clip]{lenovo-lp1738-flat-snc4.png}
\caption{Flat mode on a two-socket Xeon Max server with Quadrant
clustering (top) and SNC4 (bottom).  Blue dashed boundaries are
CPU-bearing DDR5 NUMA nodes; red dashed boundaries are separately
addressable HBM nodes.  Quadrant mode exposes one 64\,GB HBM node per
socket, whereas SNC4 exposes four 16\,GB HBM nodes per socket.  Reproduced
from Lenovo Press LP1738, Figures 7 and 11~\cite{lp1738}.}
\label{fig:lenovo-flat-topologies}
\end{figure}
\FloatBarrier

Lenovo's own STREAM experiment makes the interaction unusually visible
in \figref{fig:lenovo-hbm-bandwidth}.  The DDR-only reference used two
Xeon Platinum 8490H processors, while the HBM-mode tests used two Xeon
Max 9480 processors.  The headline ratios against DDR therefore change
the processor as well as the memory and are not a clean tier comparison.
The within-mode contrast is nevertheless informative: Flat+Quad without
binding delivered 101\% of the DDR reference, whereas explicit binding
raised it to 201\%; bound Flat+SNC4 reached 228\%.  Thus the largest
difference among the Flat bars is not merely the firmware mode but
whether the benchmark actually reaches the intended memory.  Comparing
the two bound Flat results, SNC4 adds 27 points on Lenovo's normalised
scale, or about 13\% relative to bound Quad; the guide's prose statement
of ``more than 25\%'' should therefore be read as percentage points for
this chart, not as a 25\% relative speed-up.  These are vendor STREAM
results from a specific SD650~V3 setup, not bandwidth values to inherit
for CRESCO8~\cite{lp1738}.

\begin{figure}[htbp]
\centering
% Crop the border and drop shadow embedded in the source chart.
\includegraphics[width=0.92\textwidth,trim=3.84bp 17.28bp 17.28bp 3.84bp,clip]{lenovo-lp1738-stream-bandwidth.png}
\caption{Lenovo's normalised STREAM comparison across DDR5, Cache, Flat
and HBM-only modes.  The baseline is the unbound Quadrant DDR result on
a Xeon Platinum 8490H system; the three HBM modes use Xeon Max 9480, so
the chart is not a same-processor DDR--HBM comparison.  The large
separation between bound and unbound Flat-mode bars is the important
result for this thesis: exposing HBM does not by itself place the
workload there.  Reproduced from Lenovo Press LP1738, Figure
21~\cite{lp1738}.}
\label{fig:lenovo-hbm-bandwidth}
\end{figure}
\FloatBarrier

\subsection{Cooling and instrumentation}

Direct water cooling allows two high-power processors to operate in a
dense node, but it does not guarantee a fixed frequency or energy
efficiency.  The relevant quantities are measured socket frequency,
temperature, package power and throttling state.  Product-guide claims
about facility heat recovery or cooling savings are not evidence about
inference performance.

The SD650~V3 includes XClarity management and a TI INA226-based
power/energy meter with 100\,Hz sampling capability
~\cite{lp1603}.  This establishes a platform capability, not user
access on CRESCO8.  Intel RAPL, management-controller telemetry and
whole-node power meters cover different system boundaries; they must
not be combined without first validating their scope, units, sampling
and permissions.

\FloatBarrier
\subsection{Consequences for the thesis}

Five requirements follow from the hardware, and each of them exists
because a specific thing can go wrong.

\begin{enumerate}
  \item Establish the active memory and clustering modes from the
        allocated node rather than from the product guide.  Flat mode is
        a precondition for the primary comparison, not a default.
  \item Measure sustained local HBM and DDR5 behaviour before drawing a
        Roofline or interpreting any inference result.  The published
        figures for this processor family differ by more than 15\%
        between studies, which is too much to inherit.
  \item Treat thread count as a factor rather than a constant.  It sets
        the memory-level parallelism available to the run, and therefore
        how much of the HBM bandwidth the workload can reach at all
        (\secref{sec:littles-law}).
  \item Bind threads and pages, then verify residency and remote traffic
        rather than trusting the command line.  A memory policy records
        an intention; only a page count records an outcome.
  \item Record frequency, thermal and power state throughout, so that a
        placement effect cannot be confused with a change in operating
        conditions.
\end{enumerate}

The common thread is that this processor offers several ways to obtain a
result that looks like a memory-tier effect and is not one.  Most of the
methodological weight in \chapref{chap:methodology} goes on ruling those
out.
